\section*{Hoofdstuk \ref{ch:b2:schrodinger-wave-functions} --- De
schrödingervergelijking en golffuncties}

\begin{solution}{exo:b2:schrodinger-wave-functions:1}
$h/\sqrt{2mE}$: \qty{0.123}{nm}, \qty{12.3}{pm}, \qty{3.9}{pm} (bij
\qty{100}{keV} is de
kinetische energie een vijfde van $mc^2$; de relativistische impuls geeft
\qty{3.7}{pm}); neutron \qty{0.18}{nm}; C$_{60}$ \qty{2.8}{pm}; tennisbal
$\sim 10^{-34}\,\unit{m}$. Alle behalve de bal hebben interferentie getoond.
\end{solution}

\begin{solution}{exo:b2:schrodinger-wave-functions:2}
$A = (\pi a^2)^{-1/4}$; $\langle x\rangle = 0$, $\Delta x = a/\sqrt2$;
$0.84$; $0.9995$.
\end{solution}

\begin{solution}{exo:b2:schrodinger-wave-functions:3}
Invullen geeft $E\varphi = -(\hbar^2/2m)\varphi'' + V\varphi$.
$E/h =
\qty{2.4e14}{Hz}$; \qty{4.8e14}{Hz}, \qty{620}{nm}: het rode licht
dat de overgang
uitzendt.
\end{solution}

\begin{solution}{exo:b2:schrodinger-wave-functions:4}
Elektron, \qty{1}{nm}: $\tau = \qty{1.7e-14}{s}$, \qty{60}{\micro m} na \qty{1}{ns},
\qty{60}{km} na \qty{1}{s}; \qty{1}{mm}: $\tau = \qty{17}{s}$, onveranderd
na \qty{1}{ns},
$+0.2\%$ na \qty{1}{s}; proton, \qty{1}{nm}: $\tau = \qty{3.2e-11}{s}$,
\qty{30}{nm},
\qty{30}{m}; korrel: $\tau = 600$ jaar. Voor zichtbare voorwerpen is $\tau$
astronomisch --- en hun omgeving lokaliseert ze voortdurend opnieuw.
\end{solution}

\begin{solution}{exo:b2:schrodinger-wave-functions:5}
(a) \cref{prop:b2:schrodinger-wave-functions:current}. (b) $|A|^2\hbar k/m$;
$0$; $(|A|^2 - |B|^2)\hbar k/m$. (c) De invallende en weerkaatste fluxen
worden zonder interferentieterm van elkaar afgetrokken: doorlatings- en
weerkaatsings\-coëfficiënten
kunnen uit fluxen worden gedefinieerd. (d) $\psi^*\psi'$ reëel, het
imaginaire deel ervan nul.
\end{solution}

\begin{solution}{exo:b2:schrodinger-wave-functions:6}
(a) $A = (2\pi\sigma^2)^{-1/4}$, $\langle x\rangle = 0$,
$\Delta x = \sigma$. (b) $\hbar k_0$.
(c) $\sigma \times \hbar/2\sigma = \hbar/2$. (d)
$\dd\langle x\rangle/\dd t = \int x\,
\partial_t\rho = -\int x\,\partial_xj = \int j = (\hbar/m)\operatorname{Im}\int\psi^*
\psi' = \langle p\rangle/m$.
\end{solution}

\begin{solution}{exo:b2:schrodinger-wave-functions:7}
(a) \qty{5.9e6}{m/s}, \qty{0.123}{nm}, \qty{51}{ns}. (b)
$v_\varphi = \qty{3.0e6}{m/s}$,
$v_{\text{g}} = \qty{5.9e6}{m/s}$: de laatste. (c)
$\tau = \qty{1.7}{\micro s} \gg
\qty{51}{ns}$: onveranderd. (d)
$\lambda/a = \qty{4e-7}{rad}$, \qty{0.1}{\micro m}:
stralen voldoen.
\end{solution}

\begin{solution}{exo:b2:schrodinger-wave-functions:8}
(a) $|\psi|^2$ en elke verwachtingswaarde blijven onveranderd. (b)
$\tfrac12(|\psi_1|^2
+ |\psi_2|^2) + \operatorname{Re}(\eu^{\iu\alpha}\psi_1^*\psi_2)$.
(c) De golven van
de twee spleten; hun betrekkelijke fase is $k(r_2 - r_1)$. (d) De detector
komt voor de twee wegen in verschillende toestanden terecht; de twee componenten
interfereren niet meer, en het patroon is de som van twee enkelspleet\-patronen.
\end{solution}

\begin{solution}{exo:b2:schrodinger-wave-functions:9}
(a) $\langle V'\rangle = m\omega^2\langle x\rangle$: exact. (b)
$\langle V'\rangle
= mg$: exact. (c) Wanneer het pakket smal is tegenover de
schaal waarop
$V''$ verandert. (d) Zijn golflengte bij \qty{0.1}{m/s} is $10^{-29}\,\unit{m}$; het
pakket zou zo breed als de kom moeten zijn om ertoe te doen. De thermische
golflengte van een atoom, \qty{0.1}{nm}, is zijn eigen grootte.
\end{solution}

\begin{solution}{exo:b2:schrodinger-wave-functions:10}
(a)
$\tfrac12(\varphi_1^2 + \varphi_2^2) + \varphi_1\varphi_2\cos\omega_{21}t$.
(b)
$L/2 +
\cos\omega_{21}t\int x\varphi_1\varphi_2 = L/2 - (16L/9\pi^2)\cos\omega_{21}t$.
(c)
$3h^2/8mL^2 = \qty{1.13}{eV}$, \qty{2.7e14}{Hz}, \qty{1.1}{\micro m}. (d)
Zij straalt
op $\nu_{21}$: de lijn; een zuivere \omterm{thm:b2:schrodinger-wave-functions:stationary}{stationaire toestand} heeft geen trillende
dipool en straalt niet.
\end{solution}

\begin{solution}{exo:b2:schrodinger-wave-functions:11}
(a)
$\psi \propto \int\eu^{-\sigma_0^2k^2}\eu^{\iu(kx - \hbar k^2t/2m)}\dd k$.
(b) $\alpha =
\sigma_0^2 + \iu\hbar t/2m$, $\beta = \iu x$:
$|\psi|^2 \propto \exp[-x^2/2\sigma_0^2(1 +
(\hbar t/2m\sigma_0^2)^2)]$. (c)
$\Delta v = \hbar/2m\sigma_0$, en $\Delta v\,\tau =
\sigma_0$. (d)
$\omega = ck$: geen dispersie, alle componenten met $c$.
\end{solution}

\begin{solution}{exo:b2:schrodinger-wave-functions:12}
(a) $\Gamma = \hbar/\tau$: \qty{6.6e-8}{eV}; \qty{6.6e-8}{eV}; \qty{66}{MeV}. (b)
$1/2\pi\tau = \qty{16}{MHz}$, honderd keer minder dan Doppler. (c) In
koude atomen en gevangen ionen verdwijnt de dopplerbreedte en wordt de natuurlijke
breedte gehaald --- de grondslag van frequentiestandaarden. (d)
$c\tau =
\qty{3}{m}$, de lengte van de trein.
\end{solution}

\begin{solution}{pb:b2:schrodinger-wave-functions:1}
\textbf{1.} $v = \sqrt{2eU/m} = \qty{8.4e7}{m/s}$;
$p = \qty{7.6e-23}{kg.m/s}$; $\lambda =
\qty{8.7}{pm}$.

\textbf{2.} Fouten van enkele procenten: aanvaardbaar voor schattingen
($\lambda =
\qty{8.6}{pm}$ precies).

\textbf{3.} \qty{4.8}{ns}.

\textbf{4.} $\lambda/a = \qty{3e-8}{rad}$: \qty{12}{nm} --- $10^4$ keer
kleiner dan een
pixel.

\textbf{5.} $\Delta\theta = \Delta p_y/p \sim \hbar/2ap = \lambda/4\pi a$:
dezelfde orde.

\textbf{6.} $\tau = \qty{1.6}{ms} \gg \qty{4.8}{ns}$: onveranderd.

\textbf{7.} $\dd\langle p_y\rangle/\dd t = eE$ precies voor een gelijkmatig veld:
$\langle y\rangle$ volgt de klassieke parabool; de buis is een
klassiek instrument.

\textbf{8.} $v/2 = \qty{4.2e7}{m/s}$; alleen de \omterm{prop:b2:dispersion-wave-packets:group}{groepssnelheid} is waarneembaar;
de fase van een vlakke golf draagt geen signaal (en hangt af van het willekeurige
nulpunt van de energie).

\textbf{9.} $\Delta p = \hbar/2\sigma_0$, $\Delta v = \hbar/2m\sigma_0$.

\textbf{10.} \qty{1.7e-16}{s}, \qty{1.7e-14}{s}, \qty{1.7e-8}{s}; proton
\qty{3.2e-11}{s};
virus \qty{20}{s}.

\textbf{11.} Het is in een \omterm{thm:b2:schrodinger-wave-functions:stationary}{stationaire toestand}: de potentiaal houdt het
pakket vast; $|\varphi|^2$ beweegt niet.

\textbf{12.} $v = \qty{5.9e5}{m/s}$; verdubbeling bij
$t = \sqrt3\tau = \qty{2.9e-14}{s}$:
\qty{17}{nm}, zo'n veertien golflengten.

\textbf{13.} Hun \omterm{def:b2:schrodinger-wave-functions:psi}{golffuncties} strekken zich over het hele kristal uit: een
bepaalde impuls, geen bepaalde plaats.

\textbf{14.} Oplossend vermogen
$\sim\lambda/\alpha = \qty{4}{pm}/0.01 = \qty{0.4}{nm}$: de
afwijkingen van magnetische lenzen begrenzen de opening, niet de golflengte.

\textbf{15.} Vergelijk $\lambda$ met de betrokken afmetingen (openingen, de
schaal van $V$), en de uitspreidingstijd met de waarnemingstijd.

\textbf{16.} \qty{0.38}{}, \qty{1.50}{}, \qty{3.38}{eV};
$\nu_{21} = \qty{2.7e14}{Hz}$
(\qty{1.1}{\micro m}), $\nu_{31} = \qty{7.3e14}{Hz}$ (\qty{410}{nm}).

\textbf{17.} De fasor heeft modulus één; bij $x = L/2$.

\textbf{18.} Bij $t = 0$ hoopt de dichtheid zich op de linkerhelft op, bij
$T_{21}/2$ op de rechter.

\textbf{19.} $\langle x\rangle = L/2 - (16L/9\pi^2)\cos\omega_{21}t$: amplitude
$0.18L = \qty{0.18}{nm}$.

\textbf{20.} Op $\nu_{21}$: de spectraallijn; een \omterm{thm:b2:schrodinger-wave-functions:stationary}{stationaire toestand} heeft geen
trillende dipool en is stabiel; de superpositie straalt zich omlaag
tot $\varphi_1$.

\textbf{21.} Niet bepaald: $\langle E\rangle = (E_1 + E_2)/2 = \qty{0.94}{eV}$; een
meting geeft $E_1$ of $E_2$, elk met kans een half.

\textbf{22.} Willekeurig: waar elk elektron landt; niet willekeurig: het
patroon $|\psi|^2$ dat de landingen opbouwen.

\textbf{23.} $j = |A|^2\hbar k/m$: met $|A|^2 = n$ elektronen per lengte-eenheid,
$nv$ elektronen per seconde.

\textbf{24.} Van eerste orde opdat $\psi(x,0)$ een volledige beschrijving is
die de toekomst vastlegt; complex opdat een bepaalde energie de
enkele tijdafhankelijkheid $\eu^{-\iu Et/\hbar}$ heeft en $E = p^2/2m$ eruit komt.

\textbf{25.} $\psi$ is de amplitude waarvan het kwadraat de kans is; zij
voldoet aan de lineaire vergelijking van Schrödinger; \omterm{thm:b2:schrodinger-wave-functions:stationary}{stationaire toestanden}
dragen een
fasor en vaste dichtheden; een vrij pakket beweegt met $p/m$ en spreidt zich
uit in $2m\sigma_0^2/\hbar$; Ehrenfest geeft Newton terug voor het grote.
\end{solution}
